\renewcommand{\footerfilename}{Community.tex}

\section{Building the Community IRAF/PyRAF}

IRAF/PyRAF was converted to Community support in 2016, version on the order of 2.16.

The current release, as of \today\; runs on most modern Linux and Mac machines, including
the Mac M4 processor.

\subsection{PyRAF}

Getting PyRAF to run on a M4 Mac is non-trivial. These recommendations are made
for pyraf 3.12 and later. The package absolutely requires X11. It may run inside
a Windows enviromment using a virtualization process together with XMing etc.
Mac requires XQuartz. \index{Apple!X11}

For the Mac install the \dhl{Xcode} package, avaliable for free from Apple.

\vspace{-.15cm}
\begin{enumerate}\addtolength{\itemsep}{-0.5\baselineskip}
%\setcounter{enumi}{N}
   \item   Create a python virtual environment:
  \dhl{/usr/lib/python3.10/python3 -m venv mypyraf}
   \item   \dhl{. mypyraf/bin/activate} to activate the environment. 
   \item   cd into this directory.
   \item   git clone the community pyraf here.
   \item   cd into the clone
   \item   edit/replace the pyproject.toml file.
   \item   run \dhl{pip3 install -e .}. Lots of action.
   \item   At this time there should be a very small pyraf executable in the envrioments
  bin directory.
\end{enumerate}

\subsubsection{Managing the pyproject.toml file}

The toml files are shifhting sands. They require tuning from time to
time.

Examine the \dhl{pyproject.toml} file and remove the license statement
and make sure:
\begingroup \fontsize{10pt}{10pt}
\selectfont
%%\begin{Verbatim} [commandchars=\\\{\}]
\begin{verbatim} 
[tool.setuptools_scm]

\end{verbatim}
\endgroup
is simply blank. 
